
\subsection{马斯洛夫类猜想的应用：定理 6.1.15 与 6.1.17 的证明}
\label{subsec:6-4-4}
\liftIndex{I022}

\noindent\textbf{定理 6.1.15 的证明。}
由定理~\ref{thm:3-1-11}
及定理~\ref{thm:3-8-50}，
障碍类落在
\[
H^{2-\mu_L(\beta_k)}(L;\mathbb Q)
\big/
\operatorname{Im}\bigl(
H^{2-\mu_L(\beta_k)}(M;\mathbb Q)
\longrightarrow
H^{2-\mu_L(\beta_k)}(L;\mathbb Q)
\bigr)
\]
中。
由于 $\mu_L=0$，
且按假设
$H^2(M;\mathbb Q)\to H^2(L;\mathbb Q)$
满射，所有障碍类均消失。
故 $L$ 在 bulk 形变后无障碍。
再由定理~\ref{thm:6-1-9}，
对所取的
$\mathbf b\in\mathcal M_{\mathrm{def}}(L)$，
\[
HF((L,\mathbf b),(L,\mathbf b);
\Lambda_{\mathrm{nov}})\ne0.
\]
\hfill$\square$

\noindent\textbf{定理 6.1.17 的证明。}
先证明 \textup{(1)}。
假设 $L$ 在 bulk 形变后弱无障碍，
而 $\Sigma_L=0$。
则存在
$\mathbf b\in\mathcal M_{\mathrm{weak,def}}(L)$。
由 $\Sigma_L=0$
及定理~\ref{thm:6-1-9}，
\[
HF((L,\mathbf b),(L,\mathbf b);
\Lambda_{\mathrm{nov}})\ne0.
\]
这与 $L$ 可位移的假设矛盾。

再证明 \textup{(2)}。
取假设中的
$(\boldsymbol b,b)
\in\widehat{\mathcal M}_{\mathrm{weak,def}}(L)$。

\begin{lemma}{6.4.24}\label{lem:6-4-24}
\[
\pi_2\bigl(G(L;\boldsymbol b,b)\bigr)
=\pi_2\bigl(G(L)\bigr),
\]
其中 $\pi_2:\mathbb R\times2\mathbb Z\to2\mathbb Z$
为第二坐标投影，
$G(L;\boldsymbol b,b)$ 由 (6.4.21) 定义。
\end{lemma}

\begin{proof}
若写
\[
\boldsymbol b
=\sum_i e^{m_i/2}T^{\lambda_i}\boldsymbol b_i,
\qquad
b=\sum_i e^{k_i/2}T^{\lambda'_i}b_i,
\]
则
$m_i=2-\deg\boldsymbol b_i=0$，
$k_i=1-\deg b_i=0$。
引理由定义即得。
\end{proof}

\pagebreak
以下使用典范模型的微分：
\[
\delta^{\mathrm{can}}_{(\boldsymbol b,\widetilde b)}
=\mathfrak m_1^{\mathrm{can},\boldsymbol b,\widetilde b}
:H(L;\Lambda_{0,\mathrm{nov}})
\longrightarrow H(L;\Lambda_{0,\mathrm{nov}}),
\]
其中 $\widetilde b$ 是与 $b$ 对应的典范模型定界余链。
复形
$\bigl(H(L;\Lambda_{0,\mathrm{nov}}),
\delta^{\mathrm{can}}_{(\boldsymbol b,\widetilde b)}\bigr)$
的上同调是
$HF((L;\boldsymbol b,b),(L;\boldsymbol b,b))$。
由引理~\ref{lem:6-4-24}，
可写\translatornote{原书实体页 390
在典范模型微分中把 $\widetilde b$ 写成链模型的 $b$；
同一和式将 $G\subset\mathbb R\times2\mathbb Z$
的坐标倒写为 $(\mu_i,\lambda_i)$，
跨页后又改用 $k_i$ 表示第二坐标。
此处依 (6.4.21) 及前小节的典范模型记号
统一为 $(\lambda_i,k_i)$；
见疑误 A-0381、A-0384。}
\[
\delta^{\mathrm{can}}_{(\boldsymbol b,\widetilde b)}
=\sum_{(\lambda_i,k_i)\in G(L;\boldsymbol b,b)}
e^{k_i/2}T^{\lambda_i}\delta_i,
\qquad
\delta_i:H(L;\mathbb C)\longrightarrow H(L;\mathbb C).
\]

现在反设
$\Sigma_L\notin\{2,\ldots,n+1\}$，
并导出矛盾。
由此假设，
对每个 $k_i$，
或者
\begin{align}
k_i&=0,
\tag{6.4.25.1}\label{eq:6-4-25-1}\\
|k_i|&>n+1.
\tag{6.4.25.2}\label{eq:6-4-25-2}
\end{align}
在 (6.4.25.2) 情形，
$\deg\delta_i=1-k_i
\notin\{-n,\ldots,n\}$。
因为 $\delta_i$ 作用于
$n$ 维流形 $L$ 的通常上同调，
故 $\delta_i=0$。
\translatornote{原书实体页 391
在只得到 (6.4.25.2) 类分量为零时
便称其与可位移性矛盾；
真正的矛盾还需下面对
$k_i=0$ 的剩余分量证明幺元类非零，
见疑误 A-0385。}
因此只须考虑 $k_i=0$ 的项。
这时
$\operatorname{PD}([L])\in H^0(L;\mathbb C)$
不在
$\delta^{\mathrm{can}}_{(\boldsymbol b,\widetilde b)}$
的像中。
另一方面，
$\operatorname{PD}([L])$ 是幺元，
故
\[
\delta^{\mathrm{can}}_{(\boldsymbol b,\widetilde b)}
\bigl(\operatorname{PD}([L])\bigr)=0.
\]
于是得自弗洛尔上同调非零：
\[
HF((L;\boldsymbol b,b),(L;\boldsymbol b,b);
\Lambda_{\mathrm{nov}})\ne0,
\]
与 $L$ 可位移矛盾。
定理~\ref{thm:6-1-17} 得证。
\hfill$\square$

\begin{remark}{6.4.26}\label{rem:6-4-26}
定理~\ref{thm:6-1-17} 的证明
只用到 $\operatorname{PD}([L])$，
未用到 $\operatorname{PD}([pt])$。
因此，当整数或 $\mathbb Z/2\mathbb Z$ 系数的
弗洛尔上同调有定义时，
考虑
\[
\operatorname{PD}([L])
-\mathfrak h_2
\bigl(e^b\mathbin{\overline\otimes}e^b\bigr),
\]
同一证明也适用于
$\Lambda_{0,\mathrm{nov}}(\mathbb Z)$
（或
$\Lambda_{0,\mathrm{nov}}(\mathbb Z/2\mathbb Z)$）
上的情形。
见 \textbf{[FOOO09I]}。
\end{remark}
